\title{2022 年高考数学 全国甲卷-文科}
\date{}
\maketitle

\section*{一、选择题}
本题共12小题，每小题5分，共60分。在每小题给出的四个选项中，只有一项是符合题目要求的。

\begin{question}[points=5]
设集合 $A=\{-2,-1,0,1,2\}, B=\{x \mid 0 \leq x < \frac{5}{2}\}$，则 $A \cap B=$
\begin{choices}
  \item $\{0,1,2\}$
  \item $\{-2,-1,0\}$
  \item $\{0,1\}$
  \item $\{1,2\}$
\end{choices}
\begin{solution}
\textbf{答案：}A
\end{solution}
\end{question}

\begin{question}[points=5]
某社区通过公益讲座以普及社区居民的垃圾分类知识．为了解讲座效果，随机抽取10位社区居民，让他们在讲座前和讲座后各回答一份垃圾分类知识问卷，这10位社区居民在讲座前和讲座后问卷答题的正确率如下图：



则

\begin{center}
  \includegraphics[width=.38\linewidth]{question-2-1.jpg}
\end{center}
\begin{choices}
  \item A．讲座前问卷答题的正确率的中位数小于70%
  \item B．讲座后问卷答题的正确率的平均数大于85%
  \item C．讲座前问卷答题的正确率的标准差小于讲座后正确率的标准差
  \item D．讲座后问卷答题的正确率的极差大于讲座前正确率的极差
\end{choices}
\begin{solution}
\textbf{答案：}B
\end{solution}
\end{question}

\begin{question}[points=5]
若 $z=1+\mathrm{i}$，则 $|\mathrm{i}z+3\overline{z}|=$
\begin{choices}
  \item $4\sqrt{5}$
  \item $4\sqrt{2}$
  \item $2\sqrt{5}$
  \item $2\sqrt{2}$
\end{choices}
\begin{solution}
\textbf{答案：}D
\end{solution}
\end{question}

\begin{question}[points=5]
如图，网格纸上绘制的是一个多面体的三视图，网格小正方形的边长为1，则该多面体的体积为
\begin{choices}
  \item A．8
  \item B．12
  \item C．16
  \item D．20
\end{choices}
\begin{solution}
\textbf{答案：}B
\end{solution}
\end{question}

\begin{question}[points=5]
将函数 $f(x)=\sin\left(\omega x+\frac{\pi}{3}\right) (\omega>0)$ 的图像向左平移 $\frac{\pi}{2}$ 个单位长度后得到曲线 $C$，若 $C$ 关于 $y$ 轴对称，则 $\omega$ 的最小值是
\begin{choices}
  \item $\frac{1}{6}$
  \item $\frac{1}{4}$
  \item $\frac{1}{3}$
  \item $\frac{1}{2}$
\end{choices}
\begin{solution}
\textbf{答案：}C
\end{solution}
\end{question}

\begin{question}[points=5]
从分别写有1，2，3，4，5，6的6张卡片中无放回随机抽取2张，则抽到的2张卡片上的数字之积是4的倍数的概率为
\begin{choices}
  \item $\frac{1}{5}$
  \item $\frac{1}{3}$
  \item $\frac{2}{5}$
  \item $\frac{2}{3}$
\end{choices}
\begin{solution}
\textbf{答案：}C
\end{solution}
\end{question}

\begin{question}[points=5]
函数 $f(x)=\left(3^x-3^{-x}\right)\cos x$ 在区间 $\left[-\frac{\pi}{2},\frac{\pi}{2}\right]$ 的图像大致为

\begin{center}
  \includegraphics[width=.38\linewidth]{question-7-1.jpg}
\end{center}
\begin{choices}
  \item A．
  \item B．
  \item C．
  \item D．
\end{choices}
\begin{solution}
\textbf{答案：}A
\end{solution}
\end{question}

\begin{question}[points=5]
当 $x=1$ 时，函数 $f(x)=a\ln x+\frac{b}{x}$ 取得最大值 $-2$，则 $f'(2)=$
\begin{choices}
  \item $-1$
  \item $-\frac{1}{2}$
  \item $\frac{1}{2}$
  \item $1$
\end{choices}
\begin{solution}
\textbf{答案：}B
\end{solution}
\end{question}

\begin{question}[points=5]
在长方体 $ABCD-A_1B_1C_1D_1$ 中，已知 $B_1D$ 与平面 $ABCD$ 和平面 $AA_1B_1B$ 所成的角均为 $30^\circ$，则
\begin{choices}
  \item A．$AB=2AD$
  \item B．$AB$ 与平面 $AB_1C_1D$ 所成的角为 $30^\circ$
  \item C．$AC=CB_1$
  \item D．$B_1D$ 与平面 $BB_1C_1C$ 所成的角为 $45^\circ$
\end{choices}
\begin{solution}
\textbf{答案：}D
\end{solution}
\end{question}

\begin{question}[points=5]
甲、乙两个圆锥的母线长相等，侧面展开图的圆心角之和为 $2\pi$，侧面积分别为 $S_{\text{甲}}$ 和 $S_{\text{乙}}$，体积分别为 $V_{\text{甲}}$ 和 $V_{\text{乙}}$．若 $\frac{S_{\text{甲}}}{S_{\text{乙}}}=2$，则 $\frac{V_{\text{甲}}}{V_{\text{乙}}}=$
\begin{choices}
  \item $\sqrt{5}$
  \item $2\sqrt{2}$
  \item $\sqrt{10}$
  \item $\frac{5\sqrt{10}}{4}$
\end{choices}
\begin{solution}
\textbf{答案：}C
\end{solution}
\end{question}

\begin{question}[points=5]
已知椭圆 $C:\frac{x^2}{a^2}+\frac{y^2}{b^2}=1(a>b>0)$ 的离心率为 $\frac{1}{3}$，$A_1,A_2$ 分别为 $C$ 的左、右顶点，$B$ 为 $C$ 的上顶点．若 $\overrightarrow{BA_1}\cdot\overrightarrow{BA_2}=-1$，则 $C$ 的方程为
\begin{choices}
  \item $\frac{x^2}{18}+\frac{y^2}{16}=1$
  \item $\frac{x^2}{9}+\frac{y^2}{8}=1$
  \item $\frac{x^2}{3}+\frac{y^2}{2}=1$
  \item $\frac{x^2}{2}+y^2=1$
\end{choices}
\begin{solution}
\textbf{答案：}B
\end{solution}
\end{question}

\begin{question}[points=5]
已知 $9^m=10, a=10^m-11, b=8^m-9$，则
\begin{choices}
  \item $a>0>b$
  \item $a>b>0$
  \item $b>a>0$
  \item $b>0>a$
\end{choices}
\begin{solution}
\textbf{答案：}A
\end{solution}
\end{question}

\section*{二、填空题}
本题共4小题，每小题5分，共20分。

\begin{question}[points=5]
已知向量 $\vec{a}=(m,3),\vec{b}=(1,m+1)$．若 $\vec{a}\perp\vec{b}$，则 $m=$
\fillin[{$-\frac{3}{4}$}]
\begin{solution}
\textbf{答案：}$-\frac{3}{4}$
\end{solution}
\end{question}

\begin{question}[points=5]
设点 $M$ 在直线 $2x+y-1=0$ 上，点 $(3,0)$ 和 $(0,1)$ 均在 $\odot M$ 上，则 $\odot M$ 的方程为
\fillin[{$(x-1)^2+(y+1)^2=5$}]
\begin{solution}
\textbf{答案：}$(x-1)^2+(y+1)^2=5$
\end{solution}
\end{question}

\begin{question}[points=5]
记双曲线 $C:\frac{x^2}{a^2}-\frac{y^2}{b^2}=1(a>0,b>0)$ 的离心率为 $e$，写出满足条件“直线 $y=2x$ 与 $C$ 无公共点”的 $e$ 的一个值
\fillin[{$2$（满足 $1<e\leq\sqrt{5}$ 皆可）}]
\begin{solution}
\textbf{答案：}$2$（满足 $1<e\leq\sqrt{5}$ 皆可）
\end{solution}
\end{question}

\begin{question}[points=5]
已知 $\triangle ABC$ 中，点 $D$ 在边 $BC$ 上，$\angle ADB=120^\circ, AD=2, CD=2BD$．当 $\frac{AC}{AB}$ 取得最小值时，$BD=$
\fillin[{$\sqrt{3}-1$}]
\begin{solution}
\textbf{答案：}$\sqrt{3}-1$
\end{solution}
\end{question}

\section*{三、解答题}
共70分。解答应写出文字说明、证明过程或演算步骤。第17~21题为必考题，每个试题考生都必须作答。第22、23题为选考题，考生根据要求作答。

\begin{question}[points=12]
甲、乙两城之间的长途客车均由 A 和 B 两家公司运营，为了解这两家公司长途客车的运行情况，随机调查了甲、乙两城之间的500个班次，得到下面列联表：



| | 准点班次数 | 未准点班次数 |

| A | 240 | 20 |

| B | 210 | 30 |



（1）根据上表，分别估计这两家公司甲、乙两城之间的长途客车准点的概率；

（2）能否有90%的把握认为甲、乙两城之间的长途客车是否准点与客车所属公司有关？



附：$K^2=\frac{n(ad-bc)^2}{(a+b)(c+d)(a+c)(b+d)}$，



$P(K^2\ge k)$ | 0.100 | 0.050 | 0.010

--- | --- | --- | ---

$k$ | 2.706 | 3.841 | 6.635
\begin{solution}
\textbf{答案：}（1）A 公司准点概率 $\frac{12}{13}$，B 公司准点概率 $\frac{7}{8}$；（2）有90%的把握
\end{solution}
\end{question}

\begin{question}[points=12]
记 $S_n$ 为数列 $\{a_n\}$ 的前 $n$ 项和．已知 $\frac{2S_n}{n}+n=2a_n+1$．

（1）证明：$\{a_n\}$ 是等差数列；

（2）若 $a_4, a_7, a_9$ 成等比数列，求 $S_n$ 的最小值．
\begin{solution}
\textbf{答案：}（1）证明见解析；（2）$-78$
\end{solution}
\end{question}

\begin{question}[points=12]
小明同学参加综合实践活动，设计了一个封闭的包装盒，包装盒如图所示：底面 $ABCD$ 是边长为8（单位：cm）的正方形，$\triangle EAB, \triangle FBC, \triangle GCD, \triangle HDA$ 均为正三角形，且它们所在的平面都与平面 $ABCD$ 垂直．

（1）证明：$EF\parallel$ 平面 $ABCD$；

（2）求该包装盒的容积（不计包装盒材料的厚度）．

\begin{center}
  \includegraphics[width=.38\linewidth]{question-19-1.jpg}
\end{center}
\begin{solution}
\textbf{答案：}（1）证明见解析；（2）$\frac{640}{3}\sqrt{3}$
\end{solution}
\end{question}

\begin{question}[points=12]
已知函数 $f(x)=x^3-x, g(x)=x^2+a$，曲线 $y=f(x)$ 在点 $(x_1, f(x_1))$ 处的切线也是曲线 $y=g(x)$ 的切线．

（1）若 $x_1=-1$，求 $a$；

（2）求 $a$ 的取值范围．
\begin{solution}
\textbf{答案：}（1）$a=3$；（2）$[-1,+\infty)$
\end{solution}
\end{question}

\begin{question}[points=12]
设抛物线 $C:y^2=2px(p>0)$ 的焦点为 $F$，点 $D(p,0)$，过 $F$ 的直线交 $C$ 于 $M,N$ 两点．当直线 $MD$ 垂直于 $x$ 轴时，$|MF|=3$．

（1）求 $C$ 的方程；

（2）设直线 $MD, ND$ 与 $C$ 的另一个交点分别为 $A,B$，记直线 $MN, AB$ 的倾斜角分别为 $\alpha, \beta$．当 $\alpha-\beta$ 取得最大值时，求直线 $AB$ 的方程．
\begin{solution}
\textbf{答案：}（1）$y^2=4x$；（2）$x=2y+4$
\end{solution}
\end{question}

\begin{question}[points=10]
[选修4-4：坐标系与参数方程]

在直角坐标系 $xOy$ 中，曲线 $C_1$ 的参数方程为 $\begin{cases} x=\frac{2+t}{6} \\ y=\sqrt{t} \end{cases}$（$t$ 为参数），曲线 $C_2$ 的参数方程为 $\begin{cases} x=-\frac{2+s}{6} \\ y=-\sqrt{s} \end{cases}$（$s$ 为参数）．

（1）写出 $C_1$ 的普通方程；

（2）以坐标原点为极点，$x$ 轴正半轴为极轴建立极坐标系，曲线 $C_3$ 的极坐标方程为 $2\cos\theta-\sin\theta=0$，求 $C_3$ 与 $C_1$ 交点的直角坐标，及 $C_3$ 与 $C_2$ 交点的直角坐标．
\begin{solution}
\textbf{答案：}（1）$y^2=6x-2\ (y\ge0)$；
（2）$C_3$ 与 $C_1$ 的交点坐标为 $(\frac{1}{2},1),(1,2)$；$C_3$ 与 $C_2$ 的交点坐标为 $(-\frac{1}{2},-1),(-1,-2)$
\end{solution}
\end{question}

\begin{question}[points=10]
[选修4-5：不等式选讲]

已知 $a,b,c$ 均为正数，且 $a^2+b^2+4c^2=3$，证明：

（1）$a+b+2c\le3$；

（2）若 $b=2c$，则 $\frac{1}{a}+\frac{1}{c}\ge3$．
\begin{solution}
\textbf{答案：}（1）证明见解析；（2）证明见解析
\end{solution}
\end{question}

